\section{Facade}
\label{sec:pat-facade}

\problemstatement{Provide a single, simplified interface to a complex
subsystem of many classes, so common tasks become easy while the full
subsystem remains available for advanced use.}

\begin{patternbox}[When You Reach for It]
The trigger is \emph{``clients must call five classes in the right order
just to do one common thing.''} When using a subsystem requires knowing
too much of its internals, a Facade offers a clean front door for the
typical case.
\end{patternbox}

\subsection*{Understanding the pattern}

The Facade is a class that knows the subsystem's classes and orchestrates
them. It exposes a few high-level methods (\texttt{watchMovie()}) that
internally coordinate the amplifier, projector, lights and player. Clients
talk to the Facade; they are shielded from the subsystem's structure and
call order. Crucially the subsystem is not hidden -- power users can still
reach the underlying classes directly.

\begin{ideabox}[A Simple Front Door, Not a Wall]
A Facade reduces coupling between clients and a subsystem by centralising
the common orchestration in one place. It does not add behaviour or
restrict access; it simply makes the easy things easy while leaving the
hard things possible.
\end{ideabox}

\subsection*{C++ implementation}

\begin{lstlisting}
#include <bits/stdc++.h>
using namespace std;

struct Amplifier { void on() { cout << "Amp on\n"; } void setVolume(int v) { cout << "Volume " << v << "\n"; } };
struct Projector { void on() { cout << "Projector on\n"; } void wide() { cout << "Widescreen\n"; } };
struct Lights    { void dim(int p) { cout << "Lights " << p << "%\n"; } };
struct Player    { void play(const string& m) { cout << "Playing " << m << "\n"; } };

class HomeTheaterFacade {                       // the simple front door
    Amplifier amp; Projector projector; Lights lights; Player player;
public:
    void watchMovie(const string& movie) {      // one call orchestrates all
        lights.dim(10);
        projector.on();
        projector.wide();
        amp.on();
        amp.setVolume(7);
        player.play(movie);
    }
};

int main() {
    HomeTheaterFacade theater;
    theater.watchMovie("Inception");            // client does one thing
    return 0;
}
\end{lstlisting}

\begin{complexitybox}[Trade-offs]
\textbf{Pros.} Simplifies client code; decouples clients from subsystem
internals; a natural place to enforce a correct call sequence.
\textbf{Cons.} Can become a ``god object'' if it accumulates too much;
does not prevent misuse of the subsystem directly. Keep it thin -- a
coordinator, not a kitchen sink.
\end{complexitybox}

\begin{takeawaybox}
Facade wraps a complicated subsystem behind a small, task-oriented
interface, lowering the knowledge a client needs. Distinguish from
Adapter (converts one interface to another expected one) -- Facade invents
a \emph{new, simpler} interface over \emph{many} classes.
\end{takeawaybox}
